\documentclass[11pt,a4paper]{article}
\usepackage[utf8]{inputenc}
\usepackage[T1]{fontenc}
\usepackage[catalan]{babel}
\usepackage{lmodern}
\usepackage[sfdefault,lining]{FiraSans}
\usepackage[a4paper,top=2.3cm,bottom=2.2cm,left=1.8cm,right=1.8cm,headheight=1.3cm,headsep=0.4cm,footskip=1cm]{geometry}
\usepackage{amsmath,amssymb}
\allowdisplaybreaks
\usepackage{icomma}
\usepackage{graphicx}
\usepackage[table]{xcolor}
\usepackage{tikz}
\usepackage{fontawesome5}
\usepackage{ccicons}
\usepackage{fancyhdr}
\usepackage{lastpage}
\usepackage{tasks}
\usepackage{needspace}
\usepackage[most]{tcolorbox}
\usepackage{array}
\usepackage[hidelinks,pdftitle={Fitxa d'exercicis de la unitat 8: Representació gràfica de funcions. 2n de Batxillerat MII},pdfsubject={Matemàtiques II},pdfkeywords={representació de funcions, domini, simetries, periodicitat, asímptotes, branques parabòliques, monotonia, curvatura, funcions radicals, trigonomètriques, a trossos, transformacions, PAU}]{hyperref}

\definecolor{navy}{HTML}{1B2748}
\definecolor{accent}{HTML}{E8603A}
\definecolor{teal}{HTML}{1F9E90}
\definecolor{soft}{HTML}{EEF1F8}
\definecolor{muted}{HTML}{5D6685}

\setlength{\parindent}{0pt}
\setlength{\parskip}{3pt}
\settasks{label=\alph*), label-format=\bfseries\color{navy}, label-width=1.3em, item-indent=1.8em, column-sep=0.6em, after-item-skip=5pt, before-skip=3pt, after-skip=4pt}

% ---------- capçalera i peu ----------
\newcommand{\llicencia}{\href{https://creativecommons.org/licenses/by-nc-sa/4.0/deed.ca}{\ccbyncsa}\ \ Llicència Creative Commons Reconeixement-NoComercial-CompartirIgual 4.0 Internacional (CC BY-NC-SA 4.0)}
\pagestyle{fancy}
\fancyhf{}
\fancyhead[L]{\raisebox{-0.35\height}{\includegraphics[height=0.95cm]{logo-gva.png}}}
\fancyhead[C]{\small\color{muted}Fitxa d'exercicis · Unitat 8: Representació gràfica de funcions · 2n Batxillerat MII}
\fancyhead[R]{\raisebox{-0.35\height}{\includegraphics[height=0.95cm]{logo-vermellar.png}}}
\fancyfoot[L]{\scriptsize\color{muted}\llicencia}
\fancyfoot[R]{\small\color{muted}\thepage}
\renewcommand{\headrulewidth}{0.4pt}
\renewcommand{\headrule}{\color{muted!40}\hrule height \headrulewidth}
\fancypagestyle{primera}{\fancyhf{}\renewcommand{\headrulewidth}{0pt}
  \fancyfoot[L]{\scriptsize\color{muted}\llicencia}
  \fancyfoot[R]{\small\color{muted}\thepage}}

% ---------- nivells i exercicis ----------
\newcommand{\estrelles}[1]{{\color{accent}\ifcase#1\or\faStar\faStar[regular]\faStar[regular]\or\faStar\faStar\faStar[regular]\or\faStar\faStar\faStar\fi}}
\newcounter{exer}
\newcommand{\ex}[1]{\par\needspace{4\baselineskip}\medskip\refstepcounter{exer}\noindent{\color{navy}\bfseries\theexer.}\ \ }

\newcommand{\seccio}[2]{\par\needspace{10\baselineskip}\bigskip
  \begin{tcolorbox}[colback=navy,colframe=navy,coltext=white,arc=2mm,boxrule=0pt,left=3mm,right=3mm,top=1.2mm,bottom=1.2mm]
  {\large\bfseries #1\quad #2}
  \end{tcolorbox}\vspace{-2pt}}


\setlength{\emergencystretch}{3em}

% quadre de recordatori (només a la fitxa d'exercicis)
\newcommand{\recorda}[1]{\par\smallskip\begin{tcolorbox}[colback=accent!7,colframe=accent!55,arc=2mm,boxrule=0.5pt,left=2.5mm,right=2.5mm,top=1mm,bottom=1mm]\small{\color{accent}\faLightbulb[regular]}\ \textbf{Recorda:}\ #1\end{tcolorbox}}
% etiqueta dels exercicis de la PAU
\newcommand{\pau}[1]{{\color{teal}\small\bfseries\faUniversity\ PAU · #1.}\ }
\newcommand{\res}[1]{#1}\newcommand{\rest}[1]{#1}

% =====================================================================
% Capçalera FentMates de les fitxes: un color per a cada curs.
% Els colors són els de la pàgina inicial (LEVELS a assets/js/materials.js);
% el logo ve de _eines/logo/pdf/ (python3 _eines/logo/genera.py).
%
% Ús, al final del preàmbul de la unitat:
%   % =====================================================================
% Capçalera FentMates de les fitxes: un color per a cada curs.
% Els colors són els de la pàgina inicial (LEVELS a assets/js/materials.js);
% el logo ve de _eines/logo/pdf/ (python3 _eines/logo/genera.py).
%
% Ús, al final del preàmbul de la unitat:
%   % =====================================================================
% Capçalera FentMates de les fitxes: un color per a cada curs.
% Els colors són els de la pàgina inicial (LEVELS a assets/js/materials.js);
% el logo ve de _eines/logo/pdf/ (python3 _eines/logo/genera.py).
%
% Ús, al final del preàmbul de la unitat:
%   \input{../comu/fentmates.tex}
%   \fmcurs{1eso}
% i, al començament del document:
%   \capcalera{Matemàtiques}{1r ESO}{Fitxa d'exercicis · Unitat 2}{Divisibilitat}
% Opcional abans de \input: \def\fmneed{14} (espai mínim abans de cada secció).
% =====================================================================
\makeatletter
\def\fm@logodir{../../logo/pdf}
\providecommand{\fmneed}{10}
\def\fm@cursnom{}

% clau del curs, color i color del text damunt del color
\newcommand{\fm@curs}[3]{\@namedef{fm@color@#1}{#2}\@namedef{fm@text@#1}{#3}}
\fm@curs{1eso}{E6B84A}{navy}
\fm@curs{2eso}{3B82F6}{white}
\fm@curs{3eso}{FF7A50}{white}
\fm@curs{4eso-a}{A78BFA}{white}
\fm@curs{4eso-b}{6D28D9}{white}
\fm@curs{1batx-mi}{9B1C1C}{white}
\fm@curs{1batx-ccss}{EF6B6B}{white}
\fm@curs{2batx-mii}{1E6B47}{white}
\fm@curs{2batx-ccss}{2BB3A3}{white}

\newcommand{\fmcurs}[1]{%
  \@ifundefined{fm@color@#1}{\PackageError{fentmates}{Curs desconegut: #1}{}}{}%
  \edef\fm@tmp{\noexpand\definecolor{curs}{HTML}{\@nameuse{fm@color@#1}}%
               \noexpand\colorlet{curstext}{\@nameuse{fm@text@#1}}}\fm@tmp
  \colorlet{cursfosc}{curs!72!black}%
  \def\fm@logo{\fm@logodir/vertical-#1.pdf}%
  \def\fm@marca{\fm@logodir/marca-#1.pdf}}

% ---------- capçalera de la primera pàgina ----------
% #1 matèria (petit), #2 curs, #3 tipus de document i unitat, #4 títol
\newcommand{\capcalera}[4]{%
  \gdef\fm@cursnom{#2}%
  \par\noindent
  \begin{tcolorbox}[colback=curs,colframe=curs,boxrule=0pt,arc=2.2mm,width=3.5cm,height=2.15cm,
      valign=center,left=3mm,right=2mm,top=1mm,bottom=1mm,nobeforeafter,box align=center]
    \color{curstext}{\footnotesize #1}\par\vspace{1pt}{\fontsize{19}{21}\selectfont\bfseries #2\par}%
  \end{tcolorbox}\hspace{5mm}%
  \begin{minipage}[c]{\dimexpr\linewidth-3.5cm-5mm-2.9cm\relax}
    {\color{cursfosc}\small\bfseries\MakeUppercase{#3}}\par\vspace{2pt}
    {\color{navy}\fontsize{20}{23}\selectfont\bfseries #4\par}
  \end{minipage}\hfill
  \begin{minipage}[c]{2.6cm}\raggedleft\includegraphics[width=2.6cm]{\fm@logo}\end{minipage}\par
  \vspace{2pt}}

% ---------- barres de secció del color del curs ----------
\renewcommand{\seccio}[2]{\par\needspace{\fmneed\baselineskip}\bigskip
  \begin{tcolorbox}[colback=curs,colframe=curs,coltext=curstext,arc=2mm,boxrule=0pt,left=3mm,right=3mm,top=1.2mm,bottom=1.2mm]
  {\large\bfseries #1\quad #2}
  \end{tcolorbox}\vspace{-2pt}}

% ---------- estrelles de dificultat del color del curs (a Batxillerat no se'n posen) ----------
\@ifundefined{estrelles}{}{%
  \renewcommand{\estrelles}[1]{{\color{curs}\ifcase#1\or\faStar\faStar[regular]\faStar[regular]\or\faStar\faStar\faStar[regular]\or\faStar\faStar\faStar\fi}}}

% ---------- capçalera i peu de les pàgines ----------
% primera pàgina: logos institucionals al peu
\fancypagestyle{primera}{\fancyhf{}\renewcommand{\headrulewidth}{0pt}
  \fancyfoot[L]{\raisebox{-0.32\height}{\includegraphics[height=0.78cm]{logo-gva.png}}\hspace{5mm}%
    \raisebox{-0.32\height}{\includegraphics[height=0.78cm]{logo-vermellar.png}}}
  \fancyfoot[R]{\scriptsize\color{muted}\href{https://creativecommons.org/licenses/by-nc-sa/4.0/deed.ca}{\ccbyncsa}\ \ CC BY-NC-SA 4.0 · fentmates.org\qquad\small\thepage}}
% la resta: el símbol a l'esquerra i el curs a la dreta, amb el títol al centre com abans
\fancyhead[L]{\raisebox{-0.3\height}{\includegraphics[height=0.62cm]{\fm@marca}}}
\fancyhead[R]{\small\bfseries\color{cursfosc}\fm@cursnom}
\renewcommand{\headrule}{\color{curs!45}\hrule height \headrulewidth}
\makeatother

%   \fmcurs{1eso}
% i, al començament del document:
%   \capcalera{Matemàtiques}{1r ESO}{Fitxa d'exercicis · Unitat 2}{Divisibilitat}
% Opcional abans de \input: \def\fmneed{14} (espai mínim abans de cada secció).
% =====================================================================
\makeatletter
\def\fm@logodir{../../logo/pdf}
\providecommand{\fmneed}{10}
\def\fm@cursnom{}

% clau del curs, color i color del text damunt del color
\newcommand{\fm@curs}[3]{\@namedef{fm@color@#1}{#2}\@namedef{fm@text@#1}{#3}}
\fm@curs{1eso}{E6B84A}{navy}
\fm@curs{2eso}{3B82F6}{white}
\fm@curs{3eso}{FF7A50}{white}
\fm@curs{4eso-a}{A78BFA}{white}
\fm@curs{4eso-b}{6D28D9}{white}
\fm@curs{1batx-mi}{9B1C1C}{white}
\fm@curs{1batx-ccss}{EF6B6B}{white}
\fm@curs{2batx-mii}{1E6B47}{white}
\fm@curs{2batx-ccss}{2BB3A3}{white}

\newcommand{\fmcurs}[1]{%
  \@ifundefined{fm@color@#1}{\PackageError{fentmates}{Curs desconegut: #1}{}}{}%
  \edef\fm@tmp{\noexpand\definecolor{curs}{HTML}{\@nameuse{fm@color@#1}}%
               \noexpand\colorlet{curstext}{\@nameuse{fm@text@#1}}}\fm@tmp
  \colorlet{cursfosc}{curs!72!black}%
  \def\fm@logo{\fm@logodir/vertical-#1.pdf}%
  \def\fm@marca{\fm@logodir/marca-#1.pdf}}

% ---------- capçalera de la primera pàgina ----------
% #1 matèria (petit), #2 curs, #3 tipus de document i unitat, #4 títol
\newcommand{\capcalera}[4]{%
  \gdef\fm@cursnom{#2}%
  \par\noindent
  \begin{tcolorbox}[colback=curs,colframe=curs,boxrule=0pt,arc=2.2mm,width=3.5cm,height=2.15cm,
      valign=center,left=3mm,right=2mm,top=1mm,bottom=1mm,nobeforeafter,box align=center]
    \color{curstext}{\footnotesize #1}\par\vspace{1pt}{\fontsize{19}{21}\selectfont\bfseries #2\par}%
  \end{tcolorbox}\hspace{5mm}%
  \begin{minipage}[c]{\dimexpr\linewidth-3.5cm-5mm-2.9cm\relax}
    {\color{cursfosc}\small\bfseries\MakeUppercase{#3}}\par\vspace{2pt}
    {\color{navy}\fontsize{20}{23}\selectfont\bfseries #4\par}
  \end{minipage}\hfill
  \begin{minipage}[c]{2.6cm}\raggedleft\includegraphics[width=2.6cm]{\fm@logo}\end{minipage}\par
  \vspace{2pt}}

% ---------- barres de secció del color del curs ----------
\renewcommand{\seccio}[2]{\par\needspace{\fmneed\baselineskip}\bigskip
  \begin{tcolorbox}[colback=curs,colframe=curs,coltext=curstext,arc=2mm,boxrule=0pt,left=3mm,right=3mm,top=1.2mm,bottom=1.2mm]
  {\large\bfseries #1\quad #2}
  \end{tcolorbox}\vspace{-2pt}}

% ---------- estrelles de dificultat del color del curs (a Batxillerat no se'n posen) ----------
\@ifundefined{estrelles}{}{%
  \renewcommand{\estrelles}[1]{{\color{curs}\ifcase#1\or\faStar\faStar[regular]\faStar[regular]\or\faStar\faStar\faStar[regular]\or\faStar\faStar\faStar\fi}}}

% ---------- capçalera i peu de les pàgines ----------
% primera pàgina: logos institucionals al peu
\fancypagestyle{primera}{\fancyhf{}\renewcommand{\headrulewidth}{0pt}
  \fancyfoot[L]{\raisebox{-0.32\height}{\includegraphics[height=0.78cm]{logo-gva.png}}\hspace{5mm}%
    \raisebox{-0.32\height}{\includegraphics[height=0.78cm]{logo-vermellar.png}}}
  \fancyfoot[R]{\scriptsize\color{muted}\href{https://creativecommons.org/licenses/by-nc-sa/4.0/deed.ca}{\ccbyncsa}\ \ CC BY-NC-SA 4.0 · fentmates.org\qquad\small\thepage}}
% la resta: el símbol a l'esquerra i el curs a la dreta, amb el títol al centre com abans
\fancyhead[L]{\raisebox{-0.3\height}{\includegraphics[height=0.62cm]{\fm@marca}}}
\fancyhead[R]{\small\bfseries\color{cursfosc}\fm@cursnom}
\renewcommand{\headrule}{\color{curs!45}\hrule height \headrulewidth}
\makeatother

%   \fmcurs{1eso}
% i, al començament del document:
%   \capcalera{Matemàtiques}{1r ESO}{Fitxa d'exercicis · Unitat 2}{Divisibilitat}
% Opcional abans de \input: \def\fmneed{14} (espai mínim abans de cada secció).
% =====================================================================
\makeatletter
\def\fm@logodir{../../logo/pdf}
\providecommand{\fmneed}{10}
\def\fm@cursnom{}

% clau del curs, color i color del text damunt del color
\newcommand{\fm@curs}[3]{\@namedef{fm@color@#1}{#2}\@namedef{fm@text@#1}{#3}}
\fm@curs{1eso}{E6B84A}{navy}
\fm@curs{2eso}{3B82F6}{white}
\fm@curs{3eso}{FF7A50}{white}
\fm@curs{4eso-a}{A78BFA}{white}
\fm@curs{4eso-b}{6D28D9}{white}
\fm@curs{1batx-mi}{9B1C1C}{white}
\fm@curs{1batx-ccss}{EF6B6B}{white}
\fm@curs{2batx-mii}{1E6B47}{white}
\fm@curs{2batx-ccss}{2BB3A3}{white}

\newcommand{\fmcurs}[1]{%
  \@ifundefined{fm@color@#1}{\PackageError{fentmates}{Curs desconegut: #1}{}}{}%
  \edef\fm@tmp{\noexpand\definecolor{curs}{HTML}{\@nameuse{fm@color@#1}}%
               \noexpand\colorlet{curstext}{\@nameuse{fm@text@#1}}}\fm@tmp
  \colorlet{cursfosc}{curs!72!black}%
  \def\fm@logo{\fm@logodir/vertical-#1.pdf}%
  \def\fm@marca{\fm@logodir/marca-#1.pdf}}

% ---------- capçalera de la primera pàgina ----------
% #1 matèria (petit), #2 curs, #3 tipus de document i unitat, #4 títol
\newcommand{\capcalera}[4]{%
  \gdef\fm@cursnom{#2}%
  \par\noindent
  \begin{tcolorbox}[colback=curs,colframe=curs,boxrule=0pt,arc=2.2mm,width=3.5cm,height=2.15cm,
      valign=center,left=3mm,right=2mm,top=1mm,bottom=1mm,nobeforeafter,box align=center]
    \color{curstext}{\footnotesize #1}\par\vspace{1pt}{\fontsize{19}{21}\selectfont\bfseries #2\par}%
  \end{tcolorbox}\hspace{5mm}%
  \begin{minipage}[c]{\dimexpr\linewidth-3.5cm-5mm-2.9cm\relax}
    {\color{cursfosc}\small\bfseries\MakeUppercase{#3}}\par\vspace{2pt}
    {\color{navy}\fontsize{20}{23}\selectfont\bfseries #4\par}
  \end{minipage}\hfill
  \begin{minipage}[c]{2.6cm}\raggedleft\includegraphics[width=2.6cm]{\fm@logo}\end{minipage}\par
  \vspace{2pt}}

% ---------- barres de secció del color del curs ----------
\renewcommand{\seccio}[2]{\par\needspace{\fmneed\baselineskip}\bigskip
  \begin{tcolorbox}[colback=curs,colframe=curs,coltext=curstext,arc=2mm,boxrule=0pt,left=3mm,right=3mm,top=1.2mm,bottom=1.2mm]
  {\large\bfseries #1\quad #2}
  \end{tcolorbox}\vspace{-2pt}}

% ---------- estrelles de dificultat del color del curs (a Batxillerat no se'n posen) ----------
\@ifundefined{estrelles}{}{%
  \renewcommand{\estrelles}[1]{{\color{curs}\ifcase#1\or\faStar\faStar[regular]\faStar[regular]\or\faStar\faStar\faStar[regular]\or\faStar\faStar\faStar\fi}}}

% ---------- capçalera i peu de les pàgines ----------
% primera pàgina: logos institucionals al peu
\fancypagestyle{primera}{\fancyhf{}\renewcommand{\headrulewidth}{0pt}
  \fancyfoot[L]{\raisebox{-0.32\height}{\includegraphics[height=0.78cm]{logo-gva.png}}\hspace{5mm}%
    \raisebox{-0.32\height}{\includegraphics[height=0.78cm]{logo-vermellar.png}}}
  \fancyfoot[R]{\scriptsize\color{muted}\href{https://creativecommons.org/licenses/by-nc-sa/4.0/deed.ca}{\ccbyncsa}\ \ CC BY-NC-SA 4.0 · fentmates.org\qquad\small\thepage}}
% la resta: el símbol a l'esquerra i el curs a la dreta, amb el títol al centre com abans
\fancyhead[L]{\raisebox{-0.3\height}{\includegraphics[height=0.62cm]{\fm@marca}}}
\fancyhead[R]{\small\bfseries\color{cursfosc}\fm@cursnom}
\renewcommand{\headrule}{\color{curs!45}\hrule height \headrulewidth}
\makeatother

\fmcurs{2batx-mii}

\begin{document}
\thispagestyle{primera}

\capcalera{Matemàtiques II}{2n Batx}{Fitxa d'exercicis · Unitat 8}{\raggedright Representació gràfica de funcions}
\vspace{10pt}
Nom i cognoms: \rule{8.4cm}{0.4pt}\hfill Grup: \rule{1.6cm}{0.4pt}\hfill Data: \rule{2.2cm}{0.4pt}
\vspace{8pt}

\begin{tcolorbox}[colback=soft,colframe=soft,arc=2mm,left=3mm,right=3mm,top=1.5mm,bottom=1.5mm]
\small Resol els exercicis al quadern justificant cada pas. Els exercicis marcats amb {\color{teal}\bfseries\faUniversity\ PAU} són de les proves d'accés a la universitat de València i d'altres comunitats autònomes, traduïts al valencià. Pots comprovar els resultats amb les eines dels apunts interactius (dominis, asímptotes, estudi complet d'una funció pas a pas, transformacions\dots).
\end{tcolorbox}

% =====================================================
\seccio{1}{Domini d'una funció}

\ex{1} Calcula el domini de les funcions següents:\begin{tasks}(2)
  \task $f(x)=\dfrac{x+3}{x^2-5x+6}$
  \task $f(x)=\sqrt{9-x^2}$
  \task $f(x)=\sqrt{\dfrac{x-1}{x+2}}$
  \task $f(x)=\ln\left(4-x^2\right)$
  \task $f(x)=\dfrac{1}{\sqrt{x^2-2x-3}}$
  \task $f(x)=\sqrt[3]{\dfrac{x}{x^2-1}}$
  \task $f(x)=\dfrac{e^x}{e^{2x}-1}$
  \task $f(x)=\ln(\ln x)$
\end{tasks}\recorda{Denominadors $\neq0$; radicands d'índex parell $\geq0$ (i $>0$ si l'arrel és al denominador); arguments dels logaritmes $>0$. El domini és la intersecció de totes les condicions.}

\ex{2} Domini de:\begin{tasks}(2)
  \task $f(x)=\tan3x$
  \task $f(x)=\dfrac{1}{1-\cos x}$
  \task $f(x)=\arcsin\dfrac x2$
  \task $f(x)=\sqrt{\sin x}$, en $[0,2\pi]$
  \task $f(x)=\ln(\cos x)$, en $[0,2\pi]$
\end{tasks}

\ex{2} Domini de:\begin{tasks}(1)
  \task \pau{Aragó, ordinària 2023} $f(x)=\dfrac{2x-1}{\sqrt{x^2-x-2}}$
  \task \pau{Canàries, ordinària 2024} $f(x)=\dfrac{\ln(x+2)}{3x+4}$
\end{tasks}

\ex{3} Domini de:\begin{tasks}(1)
  \task $f(x)=\sqrt{\dfrac{x^2-4}{x^2-9}}$
  \task $f(x)=\log_2\left(x^2-x\right)-\sqrt{3-x}$
\end{tasks}


% =====================================================
\seccio{2}{Talls amb els eixos i signe de la funció}

\ex{1} Troba els punts de tall amb els eixos i estudia el signe de:\begin{tasks}(2)
  \task $f(x)=x^3-x^2-6x$
  \task $f(x)=\dfrac{x^2-4x+3}{x+1}$
  \task $f(x)=\left(x^2-1\right)e^x$
  \task $f(x)=\ln\left(x^2-8\right)$
\end{tasks}\recorda{Els zeros de $f$ i els punts on no és contínua separen intervals de signe constant: n'hi ha prou amb un valor de prova en cada un.}

\ex{2} \pau{Extremadura, juny 2018} Demostra que $f(x)=x^3+x-3$ talla l'eix d'abscisses en un únic punt, que és positiu, i troba un interval de longitud $\frac12$ que el continga.

\ex{2} Estudia el signe de $f(x)=\dfrac{(x-2)^2(x+1)}{x^3}$. On talla o toca l'eix OX?

\ex{3} Quantes solucions té l'equació $x=2\sin x$? Localitza la positiva entre dos decimals consecutius.


% =====================================================
\seccio{3}{Simetries i periodicitat}

\ex{1} Estudia la simetria de les funcions següents:\begin{tasks}(2)
  \task $f(x)=x^4-3x^2+1$
  \task $f(x)=\dfrac{x}{x^2-4}$
  \task $f(x)=e^{-x^2}\cos x$
  \task $f(x)=x^3+1$
  \task $f(x)=\ln\dfrac{1+x}{1-x}$
  \task $f(x)=\dfrac{\sin x\cos x}{x}$
\end{tasks}\recorda{Parella: $f(-x)=f(x)$ (simetria respecte de OY). Senar: $f(-x)=-f(x)$ (respecte de l'origen). El domini ha de ser simètric.}

\ex{1} Calcula el període de:\begin{tasks}(3)
  \task $\sin3x$
  \task $\cos\dfrac x2$
  \task $\sin x+\sin2x$
  \task $\sin^2(2x)$
  \task $\tan\dfrac{\pi x}{4}$
  \task $|\cos x|$
\end{tasks}

\ex{2} Comprova que la gràfica de $f(x)=\dfrac{x^2+1}{x-1}$ és simètrica respecte del punt en què es tallen les seues asímptotes, i que la de $g(x)=\dfrac{1}{x^2-6x+10}$ és simètrica respecte d'una recta vertical.

\ex{3} Demostra que $f(x)=\ln\left(x+\sqrt{x^2+1}\right)$ és senar.


% =====================================================
\seccio{4}{Asímptotes verticals i horitzontals}

\ex{1} Calcula les asímptotes verticals i horitzontals (amb els límits laterals):\begin{tasks}(2)
  \task $f(x)=\dfrac{2x^2-3}{x^2-1}$
  \task $f(x)=\dfrac{x+2}{x^2-4}$
  \task $f(x)=\dfrac{3e^x}{e^x+1}$
  \task $f(x)=\dfrac{\sqrt{x^2+1}}{x}$
  \task $f(x)=\dfrac{\ln x}{x-1}$
\end{tasks}\recorda{AV $x=a$ si algun límit lateral en $a$ és infinit. AH $y=k$ si $\displaystyle\lim_{x\to\pm\infty}f(x)=k$: estudia cada costat per separat. Compte: $\sqrt{x^2}=|x|$.}

\ex{2} \pau{Andalusia, ordinària 2020} Estudia les asímptotes i la monotonia de $f(x)=\dfrac{x^2-2x-3}{x^2-1}$.

\ex{2} \pau{País Basc, juny 2018} Asímptotes, monotonia i extrems de $f(x)=\dfrac{x^2}{x^2-4}$.

\ex{2} \pau{Cantàbria, setembre 2019} Domini, asímptotes, monotonia i extrems de $f(x)=\dfrac{x+4}{x^2-7x-8}$.

\ex{2} \pau{País Basc, extraordinària 2025} Asímptotes i monotonia de $f(x)=\dfrac{x}{x^2-3x-4}$, i recta tangent en $x=0$.

\ex{3} \pau{La Rioja, juny 2018} Sigui $f(x)=x\,e^{-ax}$. Estudia'n les asímptotes segons els valors de $a$ i troba $a$ perquè tinga un extrem relatiu en $x=1$. És màxim o mínim?


% =====================================================
\seccio{5}{Asímptotes obliqües i branques parabòliques}

\ex{1} Calcula totes les asímptotes de:\begin{tasks}(1)
  \task \pau{Andalusia, juny 2017} $f(x)=\dfrac{x^2}{x-1}$ (troba'n també els extrems).
  \task \pau{Astúries, extraordinària 2020} $f(x)=\dfrac{2x^3+1}{x^2}$ (i els extrems).
\end{tasks}\recorda{$m=\displaystyle\lim_{x\to\pm\infty}\frac{f(x)}{x}$ (finit i $\neq0$), $n=\displaystyle\lim_{x\to\pm\infty}\big(f(x)-mx\big)$. En les racionals, l'AO és el quocient de la divisió.}

\ex{2} \pau{Cantàbria, ordinària 2025} Sigui $f(x)=\dfrac{x^3-8}{x^2-4}$ si $x\neq2$ i $f(2)=e^k$.\begin{tasks}(1)
  \task Troba $k$ perquè sigui contínua en $x=2$.
  \task Calcula'n les asímptotes.
  \task Escriu la recta tangent en $x=1$.
\end{tasks}

\ex{2} \pau{Canàries, ordinària 2021} Sigui $f(x)=\dfrac{ax^2-2}{b-x}$.\begin{tasks}(1)
  \task Troba $a$ i $b$ perquè $f(-2)=2$ i $f$ siga contínua en $\mathbb R-\{5\}$.
  \task Per a $a=-1$, $b=-3$, calcula les asímptotes.
\end{tasks}

\ex{2} \pau{Castella i Lleó, juliol 2018} Troba $a$ i $b$ perquè $f(x)=\dfrac1x+ax+b$ tinga un mínim relatiu en $\left(\frac12,6\right)$ i, amb aquests valors, troba'n les asímptotes.

\ex{3} \pau{La Rioja, juliol 2018} Troba $a$ perquè $\displaystyle\lim_{x\to+\infty}\left(\sqrt{9x^2+ax+1}-(3x-1)\right)=2$. Quina és, aleshores, l'asímptota obliqua de $y=\sqrt{9x^2+ax+1}$ en $+\infty$?

\ex{2} Estudia les branques infinites (asímptotes o branques parabòliques) de:\begin{tasks}(2)
  \task $f(x)=\sqrt x$
  \task $f(x)=x+\sqrt x$
  \task $f(x)=x\ln x$
  \task $f(x)=\sqrt{x^2+x+1}-x$
\end{tasks}

\ex{2} Troba $b$ perquè $f(x)=\dfrac{2x^2+bx}{x+1}$ tinga l'asímptota obliqua $y=2x-1$.


% =====================================================
\seccio{6}{Monotonia, extrems i curvatura}

\ex{1} \pau{Catalunya, extraordinària 2022} Sigui $f(x)=\dfrac{9}{x^2+x-2}$. Determina'n el domini, les asímptotes, els extrems relatius i els intervals de creixement i decreixement.\recorda{Taula de signes de $f'$ amb els zeros de $f'$ i els punts que no són del domini; $f''$ per a la curvatura.}

\ex{2} \pau{Extremadura, ordinària 2021} Asímptotes, monotonia, extrems relatius i punts d'inflexió de $f(x)=e^{-x^2}$.

\ex{2} Estudia la monotonia, els extrems i la curvatura de $f(x)=x^4-4x^3$.

\ex{2} Estudia la monotonia i la curvatura de $f(x)=x^2\ln x$ i de $g(x)=\dfrac{x-1}{x^2}$.

\ex{3} Troba $a$, $b$ i $c$ perquè $f(x)=x^3+ax^2+bx+c$ tinga un punt d'inflexió en $(1,2)$ amb tangent de pendent $-3$.


% =====================================================
\seccio{7}{Representació de funcions polinòmiques i racionals}

\ex{1} Fes l'estudi complet i representa \pau{País Basc, juliol 2019} $f(x)=x^3+3x^2-2$.\recorda{Domini, simetria, talls, asímptotes (i posició), monotonia i extrems, curvatura i inflexió; després situem tota la informació i dibuixem.}

\ex{2} Fes l'estudi complet i representa \pau{Astúries, extraordinària 2024} $f(x)=\dfrac{x^2-4}{1-x}$.

\ex{2} Fes l'estudi complet i representa \pau{La Rioja, extraordinària 2020} $f(x)=\dfrac{x+3}{(x+2)^2}$ i calcula la recta tangent en el punt d'inflexió.

\ex{2} Fes l'estudi complet i representa \pau{València, extraordinària 2025} $f(x)=\dfrac{2x}{x^2+1}$ i calcula'n els extrems absoluts.

\ex{3} Fes l'estudi complet i representa $f(x)=\dfrac{x^3}{(x-1)^2}$.


% =====================================================
\seccio{8}{Representació de funcions radicals}

\ex{1} Domini, talls i asímptotes de:\begin{tasks}(2)
  \task $f(x)=\sqrt{x^2-9}$
  \task $f(x)=\dfrac{x+1}{\sqrt{x^2+3}}$
\end{tasks}\recorda{$\sqrt{x^2+bx+c}\approx|x|+\frac b2$ quan $x\to\pm\infty$: les asímptotes poden ser diferents en cada costat. Multiplica pel conjugat per a calcular $n$.}

\ex{2} Fes l'estudi complet i representa $f(x)=\dfrac{x}{\sqrt{x-1}}$.

\ex{2} Fes l'estudi complet i representa $f(x)=x-3\sqrt[3]{x}$.

\ex{2} Fes l'estudi complet i representa $f(x)=\dfrac{x+1}{\sqrt{x^2+3}}$.

\ex{3} Sigui $f(x)=x\sqrt{4-x^2}$. Estudia'n el domini, la simetria, els extrems absoluts i els punts en què la tangent és vertical, i representa-la.


% =====================================================
\seccio{9}{Representació de funcions exponencials i logarítmiques}

\ex{1} Fes l'estudi complet i representa \pau{Castella i Lleó, juny 2018} $f(x)=x\,e^{-x}$.\recorda{Domini, simetria, talls, asímptotes (i posició), monotonia i extrems, curvatura i inflexió; després situem tota la informació i dibuixem.}

\ex{2} Fes l'estudi complet i representa \pau{País Basc, extraordinària 2024} $f(x)=2x\,e^{-2x^2}$.

\ex{2} Fes l'estudi complet i representa \pau{Cantàbria, extraordinària 2022} $f(x)=\dfrac{e^x}{x}$.

\ex{2} Fes l'estudi complet i representa \pau{Catalunya, ordinària 2021} $f(x)=\dfrac{\ln x}{x}$.

\ex{2} Fes l'estudi complet i representa $f(x)=\ln\left(x^2-4\right)$.

\ex{2} \pau{Castella i Lleó, ordinària 2023} Domini, monotonia i extrems de $f(x)=x(\ln x-1)$. Com és la gràfica a prop de $x=0$?


% =====================================================
\seccio{10}{Representació de funcions trigonomètriques}

\ex{2} Fes l'estudi complet i representa $f(x)=\sin x\,(1+\cos x)$.\recorda{Si és periòdica, estudia-la en un període. Factoritza $f'$ amb les identitats ($\cos2x=2\cos^2x-1$…).}

\ex{2} Fes l'estudi complet i representa $f(x)=\sin^2x+\cos x$ en $[0,2\pi]$.

\ex{2} Estudia $f(x)=\tan x-2x$ en $\left(-\frac\pi2,\frac\pi2\right)$: asímptotes, extrems i curvatura.

\ex{2} \pau{Cantàbria, extraordinària 2020} Sigui $f(x)=\dfrac{1-\cos x}{x}$. Calcula $\displaystyle\lim_{x\to0}f(x)$ i les asímptotes de $f$.

\ex{2} Estudia la monotonia i la curvatura de $f(x)=x+\cos x$ en $[0,2\pi]$ i representa-la.


% =====================================================
\seccio{11}{Funcions a trossos, amb valor absolut i transformacions}

\ex{1} Escriu a trossos i representa:\begin{tasks}(2)
  \task $f(x)=\left|x^2-2x-3\right|$
  \task $f(x)=x^2-4|x|+3$
  \task $f(x)=e^{-|x|}$
  \task $f(x)=|x|\,(x-2)$
\end{tasks}\recorda{$|g(x)|=g(x)$ si $g(x)\geq0$ i $-g(x)$ si $g(x)<0$. Els canvis de tros solen ser punts angulosos.}

\ex{2} A partir de les gràfiques de les funcions elementals, dibuixa:\begin{tasks}(2)
  \task $y=\sqrt{x+4}-2$
  \task $y=-\sqrt{-x}$
  \task $y=2-e^{-x}$
  \task $y=\left|\ln(x+1)\right|$
\end{tasks}

\ex{2} \pau{Galícia, juny 2017} Calcula $a$ i $b$ perquè sigui derivable en $x=3$ la funció $f(x)=\begin{cases}ax^2+b&\text{si }x<3\\ \ln(x-2)&\text{si }x\geq3\end{cases}$ i representa-la.

\ex{3} \pau{Madrid, extraordinària 2022} Sigui $f(x)=\begin{cases}x&\text{si }x\leq0\\ x\ln x&\text{si }x>0\end{cases}$.\begin{tasks}(1)
  \task Estudia la continuïtat i la derivabilitat en $x=0$.
  \task Estudia'n la monotonia i els extrems relatius, i representa-la.
\end{tasks}


% =====================================================
\seccio{12}{Informació a partir de les gràfiques de f, f' i f''}

\ex{2} Sabem que $f'(x)=(x-1)^2(x+2)$. Sense calcular $f$, estudia'n la monotonia, els extrems i la curvatura, i dibuixa'n una gràfica possible.\recorda{$f'>0\Rightarrow f$ creix; canvi de signe de $f'$ $\Rightarrow$ extrem; $f'$ creix ($f''>0$) $\Rightarrow f$ convexa; extrem de $f'$ $\Rightarrow$ inflexió de $f$.}

\ex{2} \pau{Catalunya, extraordinària 2021} Calcula $a$ i $b$ perquè $g(x)=ax^3+bx^2+1$ tinga un punt d'inflexió en $x=\frac12$ i la seua derivada en aquest punt siga $-\frac32$.

\ex{2} La gràfica de la derivada d'un polinomi $p$ de grau 3 és una paràbola que talla l'eix OX en $x=-1$ i $x=2$. Si $p(0)=1$ i $p(2)=-19$, troba $p$ i digues on té els extrems i el punt d'inflexió.

\ex{2} \pau{Galícia, extraordinària 2026} Dibuixa la gràfica d'una funció $f:\mathbb R-\{2\}\to\mathbb R$ tal que $f$, $f'$ i $f''$ tinguen el mateix signe en $(-1,2)$ i $\displaystyle\lim_{x\to2^+}f(x)=-\infty$. Explica el dibuix.

\ex{3} \pau{Galícia, ordinària 2022} Dibuixa la gràfica d'una funció contínua i no negativa en $[0,3]$ tal que $f(0)=f(3)=0$, $f''>0$ en $(0,1)$, $f''<0$ en $(2,3)$ i constant en $(1,2)$. Pot ser derivable en $x=1$?

\vfill
\begin{center}\color{muted}\small\faCheckCircle[regular]\ \ Ànim! Comprova els teus estudis amb les eines dels apunts interactius.\end{center}

\end{document}
